\section{Integral de Riemann}

Sea $f\colon [a, b] \to \R$ una función acotada, con $a < b$, y sean $m = \inf f([a, b])$ y $M = \sup f([a, b])$. Una \emph{partición} de $[a, b]$ es un subconjunto finito $P \subset [a, b]$ con $a, b \in P$. Escribimos $P = \{ t_0, \dots, t_n \}$ con $a = t_0 < t_1 < \cdots < t_n = b$. Para $i \in \{1, \dots, n\}$ sean $m_i$ y $M_i$ el ínfimo y el supremo de $f$ en $[t_{i-1}, t_i]$. Las \emph{sumas inferior y superior} de $f$ respecto de $P$ son
\begin{align*}
  s(f, P) & = \sum_{i=1}^{n} m_i (t_i - t_{i-1}), \\
  S(f, P) & = \sum_{i=1}^{n} M_i (t_i - t_{i-1}).
\end{align*}
Vale que $m(b - a) \leq s(f, P) \leq S(f, P) \leq M(b - a)$.

\begin{definition}[Refinamiento]
  Si $P$ y $G$ son particiones de un mismo intervalo y $P \subset G$, decimos que $G$ es \emph{más fina} que $P$ (o que refina a $P$).
\end{definition}

\begin{lemma}
  Sea $f\colon [a, b] \to \R$ acotada y sean $P \subset G$ particiones de $[a, b]$. Entonces $s(f, P) \leq s(f, G)$ y $S(f, G) \leq S(f, P)$.
\end{lemma}

\begin{proof}
  Alcanza con probarlo cuando $G$ tiene un único punto más que $P$, pues el caso general se obtiene agregando los puntos de a uno. Sea $P = \{ t_0, \dots, t_n \}$ y $G = P \cup \{ r \}$ con $t_{i-1} < r < t_i$. Sean $m_i = \inf_{[t_{i-1}, t_i]} f$, $m' = \inf_{[t_{i-1}, r]} f$ y $m'' = \inf_{[r, t_i]} f$. Como $[t_{i-1}, r]$ y $[r, t_i]$ están contenidos en $[t_{i-1}, t_i]$, se tiene $m_i \leq m'$ y $m_i \leq m''$. Entonces
  \begin{align*}
    s(f, G) - s(f, P) & = m'(r - t_{i-1}) + m''(t_i - r) - m_i (t_i - t_{i-1})       \\
                      & = (m' - m_i)(r - t_{i-1}) + (m'' - m_i)(t_i - r) \geq 0.
  \end{align*}
  La desigualdad para sumas superiores es análoga.
\end{proof}

\begin{corollary}
  Sea $f\colon [a, b] \to \R$ acotada. Si $P$ y $G$ son particiones de $[a, b]$, entonces $s(f, P) \leq S(f, G)$.
\end{corollary}

\begin{proof}
  Como $P \cup G$ refina a $P$ y a $G$,
  \[
    s(f, P) \leq s(f, P \cup G) \leq S(f, P \cup G) \leq S(f, G). \qedhere
  \]
\end{proof}

\begin{definition}[Integral inferior]
  Si $f\colon [a, b] \to \R$ es acotada, su \emph{integral inferior} es
  \[
    \lowint_a^b f(x) \,\dd x = \sup_P s(f, P),
  \]
  donde el supremo se toma sobre todas las particiones $P$ de $[a, b]$. Es decir, $s(f, P) \leq \lowint_a^b f(x) \,\dd x$ para toda partición $P$, y dado $\e > 0$ existe $P$ con $\lowint_a^b f(x) \,\dd x \leq s(f, P) + \e$.
\end{definition}

\begin{definition}[Integral superior]
  Análogamente, la \emph{integral superior} de $f$ es
  \[
    \upint_a^b f(x) \,\dd x = \inf_P S(f, P).
  \]
\end{definition}

Por el corolario anterior, $\lowint_a^b f(x) \,\dd x \leq \upint_a^b f(x) \,\dd x$.

\begin{lemma}
  Sea $a < c < b$. Si consideramos sólo las particiones de $[a, b]$ que contienen a $c$, los valores de las integrales superior e inferior no cambian.
\end{lemma}

\begin{proof}
  El supremo sobre todas las particiones es mayor o igual que el supremo sobre las que contienen a $c$. Recíprocamente, dada una partición $P$ cualquiera, $P' = P \cup \{c\}$ la refina y contiene a $c$, y $s(f, P) \leq s(f, P')$. Luego el supremo sobre las particiones que contienen a $c$ es mayor o igual que $s(f, P)$ para toda $P$, y entonces ambos supremos coinciden. Para la integral superior se razona igual con ínfimos.
\end{proof}

\begin{theorem}
  Sean $a < c < b$ y $f\colon [a, b] \to \R$ acotada. Entonces
  \begin{align*}
    \lowint_a^b f(x) \,\dd x & = \lowint_a^c f(x) \,\dd x + \lowint_c^b f(x) \,\dd x, \\
    \upint_a^b f(x) \,\dd x  & = \upint_a^c f(x) \,\dd x + \upint_c^b f(x) \,\dd x.
  \end{align*}
\end{theorem}

\begin{proof}
  Sean $A$ y $B$ los conjuntos de sumas inferiores de $f|_{[a, c]}$ y $f|_{[c, b]}$, respectivamente. Las particiones de $[a, b]$ que contienen a $c$ son exactamente las uniones $P_1 \cup P_2$ de una partición $P_1$ de $[a, c]$ y una partición $P_2$ de $[c, b]$, y $s(f, P_1 \cup P_2) = s(f|_{[a, c]}, P_1) + s(f|_{[c, b]}, P_2)$. Por lo tanto $A + B = \{ x + y : x \in A, \ y \in B \}$ es el conjunto de las sumas inferiores de $f$ respecto de particiones de $[a, b]$ que contienen a $c$. Por el lema anterior,
  \[
    \lowint_a^b f(x) \,\dd x = \sup(A + B) = \sup A + \sup B = \lowint_a^c f(x) \,\dd x + \lowint_c^b f(x) \,\dd x.
  \]
  Para las sumas superiores el razonamiento es análogo, usando ínfimos.
\end{proof}

\begin{prop}
  Sean $f, g\colon [a, b] \to \R$ acotadas. Entonces
  \[
    \lowint_a^b f \,\dd x + \lowint_a^b g \,\dd x \leq \lowint_a^b (f + g) \,\dd x \leq \upint_a^b (f + g) \,\dd x \leq \upint_a^b f \,\dd x + \upint_a^b g \,\dd x.
  \]
\end{prop}

\begin{proof}
  Sea $P$ una partición. En cada $[t_{i-1}, t_i]$ valen $m_i(f + g) \geq m_i(f) + m_i(g)$ y $M_i(f + g) \leq M_i(f) + M_i(g)$, de donde
  \[
    s(f, P) + s(g, P) \leq s(f + g, P), \qquad S(f + g, P) \leq S(f, P) + S(g, P).
  \]
  Dadas dos particiones $P$ y $G$, usando el refinamiento $P \cup G$,
  \[
    s(f, P) + s(g, G) \leq s(f, P \cup G) + s(g, P \cup G) \leq s(f + g, P \cup G) \leq \lowint_a^b (f + g) \,\dd x.
  \]
  Tomando supremo sobre $P$ y sobre $G$ se obtiene la primera desigualdad. La tercera se prueba de la misma forma con sumas superiores e ínfimos, y la del medio ya fue observada.
\end{proof}

\begin{prop}
  Sea $f\colon [a, b] \to \R$ acotada y $c \in \R$.
  \begin{itemize}
    \item Si $c \geq 0$, $\lowint_a^b c f \,\dd x = c \lowint_a^b f \,\dd x$ y $\upint_a^b c f \,\dd x = c \upint_a^b f \,\dd x$.
    \item Si $c < 0$, $\lowint_a^b c f \,\dd x = c \upint_a^b f \,\dd x$ y $\upint_a^b c f \,\dd x = c \lowint_a^b f \,\dd x$.
  \end{itemize}
\end{prop}

\begin{proof}
  Si $c \geq 0$, $m_i(cf) = c\, m_i(f)$ y $M_i(cf) = c\, M_i(f)$, y si $c < 0$, $m_i(cf) = c\, M_i(f)$ y $M_i(cf) = c\, m_i(f)$, por las propiedades del supremo y del ínfimo. Luego $s(cf, P) = c\, s(f, P)$ y $S(cf, P) = c\, S(f, P)$ si $c \geq 0$, y $s(cf, P) = c\, S(f, P)$ y $S(cf, P) = c\, s(f, P)$ si $c < 0$. Tomando supremo e ínfimo sobre las particiones (que invierte el orden cuando $c < 0$) se obtiene el resultado.
\end{proof}

\begin{prop}
  Sean $f, g\colon [a, b] \to \R$ acotadas con $f(x) \leq g(x)$ para todo $x \in [a, b]$. Entonces
  \[
    \lowint_a^b f \,\dd x \leq \lowint_a^b g \,\dd x \qquad \text{y} \qquad \upint_a^b f \,\dd x \leq \upint_a^b g \,\dd x.
  \]
\end{prop}

\begin{proof}
  Si $f \leq g$, entonces $m_i(f) \leq m_i(g)$ y $M_i(f) \leq M_i(g)$ para toda partición, y por lo tanto $s(f, P) \leq s(g, P)$ y $S(f, P) \leq S(g, P)$. Tomando supremo e ínfimo se obtienen las desigualdades para las integrales.
\end{proof}

\begin{definition}[Función integrable]
  Una función acotada $f\colon [a, b] \to \R$ es \emph{integrable} si $\lowint_a^b f(x) \,\dd x = \upint_a^b f(x) \,\dd x$. En ese caso, el valor común se denota $\int_a^b f(x) \,\dd x$.
\end{definition}

\begin{eg}
  La función $f\colon [a, b] \to \R$ dada por $f(x) = 1$ si $x \in \Q$ y $f(x) = 0$ si $x \notin \Q$ no es integrable.

  En efecto, como cada intervalo $[t_{i-1}, t_i]$ contiene racionales e irracionales, $m_i = 0$ y $M_i = 1$. Luego $s(f, P) = 0$ y $S(f, P) = b - a$ para toda partición $P$, y entonces $\lowint_a^b f \,\dd x = 0 \neq b - a = \upint_a^b f \,\dd x$.
\end{eg}

\begin{note}
  Sea $f\colon [a, b] \to \R$ acotada, y sean $\Sigma_{-}$ y $\Sigma_{+}$ los conjuntos de sumas inferiores y superiores de $f$. Que $f$ sea integrable significa que $\sup \Sigma_{-} = \inf \Sigma_{+}$. Esto ocurre si y sólo si para todo $\e > 0$ existen $s \in \Sigma_{-}$ y $S \in \Sigma_{+}$ con $S - s < \e$.
\end{note}

\begin{proof}
  ($\Rightarrow$) Si $I$ es el valor común, existen $s \in \Sigma_{-}$ y $S \in \Sigma_{+}$ con $s > I - \frac{\e}{2}$ y $S < I + \frac{\e}{2}$, de modo que $S - s < \e$.

  ($\Leftarrow$) Siempre $s \leq \sup \Sigma_{-} \leq \inf \Sigma_{+} \leq S$, de modo que $0 \leq \inf \Sigma_{+} - \sup \Sigma_{-} \leq S - s < \e$ para todo $\e > 0$, y entonces $\sup \Sigma_{-} = \inf \Sigma_{+}$.
\end{proof}

\begin{definition}[Oscilación]
  Dada una función acotada $f\colon [a, b] \to \R$ y un conjunto no vacío $X \subset [a, b]$, la \emph{oscilación} de $f$ en $X$ es $\omega(f, X) = \sup f(X) - \inf f(X)$.
\end{definition}

\begin{note}
  $\omega(f, X) = \sup \{ |f(x) - f(y)| : x, y \in X \}$.
\end{note}

\begin{proof}
  Para $x, y \in X$ vale $f(x) - f(y) \leq \sup f(X) - \inf f(X)$, y lo mismo intercambiando $x$ e $y$; luego $|f(x) - f(y)| \leq \omega(f, X)$. Recíprocamente, dado $\e > 0$ existen $x, y \in X$ con $f(x) > \sup f(X) - \frac{\e}{2}$ y $f(y) < \inf f(X) + \frac{\e}{2}$, y entonces $|f(x) - f(y)| > \omega(f, X) - \e$.
\end{proof}

\begin{theorem}
  Sea $f\colon [a, b] \to \R$ acotada. Son equivalentes:
  \begin{enumerate}[label=(\arabic*)]
    \item $f$ es integrable.
    \item Para todo $\e > 0$ existen particiones $P$ y $G$ de $[a, b]$ tales que $S(f, G) - s(f, P) < \e$.
    \item Para todo $\e > 0$ existe una partición $P$ de $[a, b]$ tal que $S(f, P) - s(f, P) < \e$.
    \item Para todo $\e > 0$ existe una partición $P = \{ t_0, \dots, t_n \}$ de $[a, b]$ tal que $\sum_{i=1}^{n} \omega_i (t_i - t_{i-1}) < \e$, donde $\omega_i = \omega(f, [t_{i-1}, t_i])$.
  \end{enumerate}
\end{theorem}

\begin{proof}
  (1) $\Leftrightarrow$ (2) Es la nota anterior, pues $s(f, P)$ recorre las sumas inferiores y $S(f, G)$ las superiores.

  (3) $\Leftrightarrow$ (4) Vale porque $\sum_{i=1}^{n} \omega_i (t_i - t_{i-1}) = S(f, P) - s(f, P)$.

  (3) $\Rightarrow$ (2) Es trivial, tomando $G = P$.

  (2) $\Rightarrow$ (3) Sean $P$ y $G$ con $S(f, G) - s(f, P) < \e$ y sea $P_0 = P \cup G$. Entonces
  \[
    s(f, P) \leq s(f, P_0) \leq S(f, P_0) \leq S(f, G),
  \]
  de donde $S(f, P_0) - s(f, P_0) \leq S(f, G) - s(f, P) < \e$.
\end{proof}

\section{Continuidad e integración}

\begin{theorem}
  Sea $f\colon [a, b] \to \R$ acotada y $a < c < b$. Entonces $f$ es integrable en $[a, b]$ si y sólo si $f|_{[a, c]}$ y $f|_{[c, b]}$ son integrables, y en ese caso
  \[
    \int_a^b f(x) \,\dd x = \int_a^c f(x) \,\dd x + \int_c^b f(x) \,\dd x.
  \]
\end{theorem}

\begin{proof}
  Por el teorema de aditividad de las integrales inferior y superior,
  \[
    \upint_a^b f \,\dd x - \lowint_a^b f \,\dd x = \left( \upint_a^c f \,\dd x - \lowint_a^c f \,\dd x \right) + \left( \upint_c^b f \,\dd x - \lowint_c^b f \,\dd x \right).
  \]
  Los dos términos entre paréntesis son no negativos, de modo que el miembro izquierdo es cero si y sólo si ambos lo son. Esto prueba la equivalencia, y la fórmula resulta de nuevo de la aditividad.
\end{proof}

\begin{theorem}
  Sean $f, g\colon [a, b] \to \R$ integrables y $c \in \R$.
  \begin{enumerate}[label=(\arabic*)]
    \item $c f$ es integrable y $\int_a^b c f(x) \,\dd x = c \int_a^b f(x) \,\dd x$.
    \item $f + g$ es integrable y $\int_a^b \bigl( f(x) + g(x) \bigr) \dd x = \int_a^b f(x) \,\dd x + \int_a^b g(x) \,\dd x$.
    \item Si $f(x) \leq g(x)$ para todo $x \in [a, b]$, entonces $\int_a^b f(x) \,\dd x \leq \int_a^b g(x) \,\dd x$.
    \item $|f|$ es integrable y $\left| \int_a^b f(x) \,\dd x \right| \leq \int_a^b |f(x)| \,\dd x$.
    \item $f g$ es integrable.
  \end{enumerate}
\end{theorem}

\begin{proof}
  (1), (2) y (3) se deducen de las proposiciones anteriores sobre integrales inferiores y superiores. Para (2), por ejemplo,
  \[
    \int_a^b f + \int_a^b g \leq \lowint_a^b (f + g) \leq \upint_a^b (f + g) \leq \int_a^b f + \int_a^b g,
  \]
  de modo que todas las desigualdades son igualdades. Para (1) se usa la proposición del producto por una constante, y para (3) la monotonía.

  (4) Para cualesquiera $x, y$ se tiene $\bigl| |f(x)| - |f(y)| \bigr| \leq |f(x) - f(y)|$. Por la nota sobre oscilaciones, $\omega(|f|, X) \leq \omega(f, X)$ para todo $X \subset [a, b]$ no vacío. En particular, para toda partición $P$, $S(|f|, P) - s(|f|, P) \leq S(f, P) - s(f, P)$, y por el teorema de caracterización $|f|$ es integrable. Como $-|f| \leq f \leq |f|$, por (1) y (3),
  \[
    -\int_a^b |f(x)| \,\dd x \leq \int_a^b f(x) \,\dd x \leq \int_a^b |f(x)| \,\dd x,
  \]
  es decir, $\left| \int_a^b f \,\dd x \right| \leq \int_a^b |f| \,\dd x$.

  (5) Sea $k > 0$ con $|f(x)| \leq k$ y $|g(x)| \leq k$ para todo $x \in [a, b]$. Si $x, y \in [t_{i-1}, t_i]$,
  \begin{align*}
    |f(x) g(x) - f(y) g(y)| & = |f(x) \bigl( g(x) - g(y) \bigr) + g(y) \bigl( f(x) - f(y) \bigr)|          \\
                            & \leq k\, \omega_i(g) + k\, \omega_i(f).
  \end{align*}
  Luego $\omega_i(fg) \leq k \bigl( \omega_i(f) + \omega_i(g) \bigr)$. Dado $\e > 0$, existen particiones $P_1$ y $P_2$ con $S(f, P_1) - s(f, P_1) < \frac{\e}{2k}$ y $S(g, P_2) - s(g, P_2) < \frac{\e}{2k}$. Para $P = P_1 \cup P_2$, que refina a ambas, las mismas desigualdades valen, y entonces
  \[
    S(fg, P) - s(fg, P) \leq k \bigl( S(f, P) - s(f, P) \bigr) + k \bigl( S(g, P) - s(g, P) \bigr) < \e. \qedhere
  \]
\end{proof}

\begin{theorem}
  Toda función continua $f\colon [a, b] \to \R$ es integrable.
\end{theorem}

\begin{proof}
  Como $[a, b]$ es compacto, $f$ es acotada y uniformemente continua. Dado $\e > 0$ existe $\delta > 0$ tal que $|x - y| < \delta \Rightarrow |f(x) - f(y)| < \frac{\e}{2(b - a)}$. Sea $n \in \N$ con $\frac{b - a}{n} < \delta$ y sea $P$ la partición $t_i = a + i \frac{b - a}{n}$, con $i \in \{0, \dots, n\}$. Si $x, y \in [t_{i-1}, t_i]$, entonces $|x - y| < \delta$ y $|f(x) - f(y)| < \frac{\e}{2(b - a)}$, de donde $\omega_i \leq \frac{\e}{2(b - a)}$. Luego
  \[
    \sum_{i=1}^{n} \omega_i (t_i - t_{i-1}) \leq \frac{\e}{2(b - a)} \sum_{i=1}^{n} (t_i - t_{i-1}) = \frac{\e}{2} < \e,
  \]
  y $f$ es integrable.
\end{proof}

\begin{theorem}
  Sea $f\colon [a, b] \to \R$ acotada. Si $f|_{[a, c]}$ es integrable para todo $c \in (a, b)$, entonces $f$ es integrable. Análogamente, si $f|_{[c, b]}$ es integrable para todo $c \in (a, b)$, entonces $f$ es integrable.
\end{theorem}

\begin{proof}
  Probamos la primera afirmación; la segunda es simétrica. Sea $k > 0$ con $|f(x)| \leq k$ para todo $x \in [a, b]$, y sea $\e > 0$. Elegimos $c \in (a, b)$ con $k (b - c) < \frac{\e}{4}$. Como $f|_{[a, c]}$ es integrable, existe una partición $\{ t_0, \dots, t_n \}$ de $[a, c]$ tal que $\sum_{i=1}^{n} \omega_i (t_i - t_{i-1}) < \frac{\e}{2}$. Agregando $t_{n+1} = b$ se obtiene una partición de $[a, b]$ y $\omega_{n+1} \leq 2k$, de modo que $\omega_{n+1} (t_{n+1} - t_n) \leq 2k (b - c) < \frac{\e}{2}$. Luego $\sum_{i=1}^{n+1} \omega_i (t_i - t_{i-1}) < \e$.
\end{proof}

\begin{corollary}
  Sea $f\colon [a, b] \to \R$ acotada. Si $f|_{[c, d]}$ es integrable para todos $c, d$ con $a < c < d < b$, entonces $f$ es integrable.
\end{corollary}

\begin{proof}
  Sea $m \in (a, b)$. Para cada $c \in (a, m)$, $f|_{[c, m]}$ es integrable, de modo que $f|_{[a, m]}$ es integrable por la segunda parte del teorema anterior. Análogamente, $f|_{[m, d]}$ es integrable para todo $d \in (m, b)$, y por la primera parte $f|_{[m, b]}$ es integrable. Por aditividad, $f$ es integrable en $[a, b]$.
\end{proof}

\begin{corollary}
  Toda función acotada $f\colon [a, b] \to \R$ con finitas discontinuidades es integrable.
\end{corollary}

\begin{proof}
  Sean $a = s_0 < s_1 < \cdots < s_r = b$ tales que todos los puntos de discontinuidad de $f$ pertenecen a $\{ s_0, \dots, s_r \}$. En cada $[s_{j-1}, s_j]$, la restricción $f|_{[c, d]}$ es continua, y por lo tanto integrable, para $s_{j-1} < c < d < s_j$. Por el corolario anterior $f|_{[s_{j-1}, s_j]}$ es integrable, y aplicando $r - 1$ veces la aditividad, $f$ es integrable en $[a, b]$.
\end{proof}

\begin{eg}
  Sea $f\colon [a, b] \to \R$ dada por $f(x) = 0$ si $x$ es irracional o $x = 0$, y $f(x) = \frac{1}{q}$ si $x = \frac{p}{q}$ con $p \in \Z$, $q \in \N$ y $\mathrm{mcd}(p, q) = 1$. Entonces $f$ es discontinua en todos los racionales no nulos de $[a, b]$ (un conjunto infinito), pero es integrable y $\int_a^b f(x) \,\dd x = 0$.

  Dado $\e > 0$, sea $\eta = \frac{\e}{b - a}$ y $F = \{ x \in [a, b] : f(x) \geq \eta \}$. El conjunto $F$ es finito: si $x = \frac{p}{q} \in F$ entonces $q \leq \frac{1}{\eta}$, y para cada $q$ hay finitos $p$ con $\frac{p}{q} \in [a, b]$. Sea $N = |F|$ y tomemos la partición uniforme de $[a, b]$ en $n$ intervalos con $\frac{2N (b - a)}{n} < \e$. Cada punto de $F$ pertenece a lo sumo a dos intervalos, de modo que los intervalos que intersecan a $F$ tienen longitud total menor que $\e$, y en ellos $M_i \leq 1$. En los demás intervalos todos los valores de $f$ son menores que $\eta$, luego $M_i \leq \eta$. Entonces
  \[
    S(f, P) = \sum M_i (t_i - t_{i-1}) < \e + \eta (b - a) = 2\e.
  \]
  Por lo tanto $\upint_a^b f \,\dd x = 0$. Como $f \geq 0$, $0 \leq \lowint_a^b f \,\dd x \leq \upint_a^b f \,\dd x = 0$, y entonces $f$ es integrable con integral $0$.
\end{eg}
